\SourceLeaf{41}%
% A544-P0041-B02
\[
\begin{aligned}
W&=2^{\frac12n}\rho^{\frac14n}\cos\frac12\omega\cos\omega\,\frac32\omega\\&\qquad\dotsfour\cos(\frac14n-\frac12)\omega
\end{aligned}
\]

% A544-P0041-B03
Cosinus in hoc productum ingredientes manifesto omnes positivi sunt, factor $\rho^{\frac14n}$ autem fit $=\cos45^0+i\sin45^0=(1+i)\surd\frac12$.\quad Hinc colligimus, $W$ esse productum ex $1+i$ in quantitatem realem positivam, unde necessario esse debebit

% A544-P0041-B04
\[
\begin{aligned}
W=(1+i)\surd n,\quad T=+\surd n,\quad U=+\surd n
\end{aligned}
\]

% A544-P0041-B05
\Article{19}

% A544-P0041-B06
Operae pretium erit, omnes summationes hactenus evolutas, hic in unum conspectum colligere.\quad Generaliter scilicet est

% A544-P0041-B07
\[
\begin{array}{c|c|c}
T=&U=&\text{prout }n\text{ est formae}\\\hline
\pm\surd n&\pm\surd n&4\mu\\
\pm\surd n&0&4\mu+1\\
0&0&4\mu+2\\
0&\pm\surd n&4\mu+3
\end{array}
\]

% A544-P0041-B08
et in casu eo, ubi $k$ supponitur $=1$, quantitati radicali signum positivum tribui debet.\quad Omni itaque iam rigore ea, quae pro valoribus primis ipsius $n$ in art. 3 per inductionem animadverteramus, demonstrata sunt, nihilque superest, nisi ut signa pro valoribus quibuscunque ipsius $k$ in omnibus casibus determinare doceamus.\quad Sed antequam hoc negotium in omni generalitate aggredi liceat, primo casus eos, ubi $n$ est numerus primus vel numeri primi potestas, propius considerare oportebit.

% A544-P0041-B09
\Article{20}

% A544-P0041-B10
Sit primo $n$ numerus primus impar, patetque per ea, quae in art. 10 exposuimus, esse $W=1+2\Sigma r^a=1+2\Sigma R^{ak}$, si statuatur $R=\cos\omega+i\sin\omega$, denotante $a$ ut illic indefinite omnia residua quadratica ipsius $n$ inter $1$ et $n-1$ contenta.\quad Quodsi quoque per $b$ indefinite omnia non-residua quadratica inter eosdem limites exprimimus, nullo negotio perspicitur, omnes numeros $ak$ congruos fieri secundum modulum $n$ vel omnibus $a$ vel omnibus $b$ (nullo ordinis respectu habito), prout $k$ vel residuum sit vel non-residuum.\quad Quamobrem in casu priori erit

\SourceLeaf{42}%
% A544-P0042-B02
\[
\begin{aligned}
W=1+2\Sigma R^a=1+R+R^4+R^9+\text{etc.}+R^{(n-1)^2}
\end{aligned}
\]

% A544-P0042-B03
adeoque $W=+\surd n$, si $n$ est formae $4\mu+1$, atque $W=+i\surd n$, si $n$ est formae $4\mu+3$.

% A544-P0042-B04
Contra in casu altero, ubi $k$ est non-residuum ipsius $n$, erit

% A544-P0042-B05
\[
\begin{aligned}
W=1+2\Sigma R^b
\end{aligned}
\]

% A544-P0042-B06
Hinc quum manifesto omnes $a$, $b$ complexum integrum numerorum $1$, $2$, $3$ $\ldots$ expleant, adeoque sit

% A544-P0042-B07
\[
\begin{aligned}
\Sigma R^a+\Sigma R^b=R+R^2+R^3+\text{etc.}+R^{n-1}=-1
\end{aligned}
\]

% A544-P0042-B08
fiet

% A544-P0042-B09
\[
\begin{aligned}
W=-1-2\Sigma R^a=-(1+R+R^4+R^9+\text{etc.}+R^{(n-1)^2})
\end{aligned}
\]

% A544-P0042-B10
adeoque $W=-\surd n$, si $n$ est formae $4\mu+1$, atque $W=-i\surd n$, si $n$ est formae $4\mu+3$.

% A544-P0042-B11
Hinc itaque colligitur

% A544-P0042-B12
\emph{primo}, si $n$ est formae $4\mu+1$, atque $k$ residuum quadraticum ipsius $n$,

% A544-P0042-B13
\[
\begin{aligned}
T=+\surd n,\quad U=0
\end{aligned}
\]

% A544-P0042-B14
\emph{secundo}, si $n$ est formae $4\mu+1$, atque $k$ non-residuum ipsius $n$,

% A544-P0042-B15
\[
\begin{aligned}
T=-\surd n,\quad U=0
\end{aligned}
\]

% A544-P0042-B16
\emph{tertio}, si $n$ est formae $4\mu+3$, atque $k$ residuum ipsius $n$,

% A544-P0042-B17
\[
\begin{aligned}
T=0,\quad U=+\surd n
\end{aligned}
\]

% A544-P0042-B18
\emph{quarto}, si $n$ est formae $4\mu+3$, atque $k$ non-residuum ipsius $n$

% A544-P0042-B19
\[
\begin{aligned}
T=0,\quad U=-\surd n
\end{aligned}
\]

% A544-P0042-B20
\Article{21}

% A544-P0042-B21
Sit secundo $n$ quadratum altiorve potestas numeri primi imparis $p$, statuaturque $n=p^{2\varkappa}q$, ita ut sit $q$ vel $=1$ vel $=p$.\quad Hic ante omnia observare convenit, si $\lambda$ sit integer quicunque per $p^\varkappa$ non divisibilis, fieri

\SourceLeaf{43}%
% A544-P0043-B02
\[
\begin{aligned}
&r^{\lambda\lambda}+r^{(\lambda+p^\varkappa q)^2}+r^{(\lambda+2p^\varkappa q)^2}+r^{(\lambda+3p^\varkappa q)^2}\\&\qquad+\text{etc.}+r^{(\lambda+n-p^\varkappa q)^{3}}\\
&\quad=r^{\lambda\lambda}\{1+r^{2\lambda p^\varkappa q}+r^{4\lambda p^\varkappa q}\\&\qquad+r^{6\lambda p^\varkappa q}+\text{etc.}+r^{2\lambda(n-p^\varkappa q)}\}\\&\quad=\frac{r^{\lambda\lambda}(1-r^{2\lambda n})}{1-r^{2\lambda p^\varkappa q}}=0
\end{aligned}
\]

% A544-P0043-B03
Hinc facile perspicietur, fieri

% A544-P0043-B04
\[
\begin{aligned}
W=1+r^{p^{2\varkappa}}+r^{4p^{2\varkappa}}+r^{9p^{2\varkappa}}+\text{etc.}+r^{(n-p^\varkappa)^2}
\end{aligned}
\]

% A544-P0043-B05
Termini enim reliqui progressionis

% A544-P0043-B06
\[
\begin{aligned}
1+r+r^4+r^9+\text{etc.}+r^{(n-1)^2}
\end{aligned}
\]

% A544-P0043-B07
distribui poterunt in $(p^\varkappa-1)q$ progressiones partiales, quae singulae sint $p^\varkappa$ terminorum, et per transformationem modo traditam summas evanescentes conficiant.

% A544-P0043-B08
Hinc colligitur, in casu eo, ubi fit $q=1$, sive ubi $n$ est potestas numeri primi cum exponente pari, fieri

% A544-P0043-B09
\[
\begin{aligned}
W=p^\varkappa=+\surd n,\;\text{adeoque}\;T=+\surd n,\quad U=0
\end{aligned}
\]

% A544-P0043-B10
Contra in casu eo, ubi $q=p$, sive ubi $n$ est potestas numeri primi cum exponente impari, statuemus $r^{p^{2\varkappa}}=\rho$, unde $\rho$ erit radix propria aequationis $x^p-1=0$, et quidem $\rho=\cos\frac{k}{p}360^0+i\sin\frac{k}{p}360^0$, ac dein

% A544-P0043-B11
\[
\begin{aligned}
W&=1+\rho+\rho^4+\rho^9+\text{etc.}+\rho^{(p^{\varkappa+1}-1)^2}\\&=p^\varkappa(1+\rho+\rho^4+\rho^9+\text{etc.}+\rho^{(p-1)^2})
\end{aligned}
\]

% A544-P0043-B12
Sed summa seriei $1+\rho+\rho^4+\rho^9+\text{etc.}+\rho^{(p-1)^2}$ per art. praec. determinatur, unde sponte concluditur, fieri

% A544-P0043-B13
\[
\begin{aligned}
W&=\pm\surd n=T,\quad\text{si fuerit }p\text{ formae }4\mu+1\\
W&=\pm i\surd n=iU,\quad\text{si fuerit }p\text{ formae }4\mu+3
\end{aligned}
\]

% A544-P0043-B14
signo positivo vel negativo valente, prout $k$ fuerit residuum vel non-residuum ipsius $p$.

% A544-P0043-B15
\Article{22}

% A544-P0043-B16
Facile quoque ex iis, quae in artt. 20. et 21 exposita sunt, derivatur propositio sequens, quae infra usum notabilem nobis praestabit.\quad Statuatur

% A544-P0043-B17
\[
\begin{aligned}
W'=1+r^h+r^{4h}+r^{9h}+\text{etc.}+r^{h(n-1)^2}
\end{aligned}
\]

\SourceLeaf{44}%
% A544-P0044-B02
denotante $h$ integrum quemcunque per $p$ non divisibilem, eritque in casu eo, ubi $n=p$, vel ubi $n$ est potestas ipsius $p$ cum exponente impari,

% A544-P0044-B03
\[
\begin{aligned}
W'&=W,\;\text{si fuerit }h\text{ residuum quadraticum ipsius }p\\
W'&=-W,\;\text{si fuerit }h\text{ non-residuum quadraticum ipsius }p
\end{aligned}
\]

% A544-P0044-B04
Patet enim, $W'$ oriri ex $W$, si pro $k$ substituatur $kh$; in casu priori autem $k$ et $kh$ similes erunt, in posteriori dissimiles, quatenus sunt residua vel non-residua ipsius $p$.

% A544-P0044-B05
In casu eo autem, ubi $n$ est potestas ipsius $p$ cum exponente pari, manifesto fit $W'=+\surd n$, adeoque semper $W'=W$.

% A544-P0044-B06
\Article{23}

% A544-P0044-B07
In artt. 20. 21. 22 consideravimus numeros primos impares, taliumque potestates: superest itaque casus, ubi $n$ est potestas binarii.

% A544-P0044-B08
Pro $n=2$ manifesto fit $W=1+r=0$.

% A544-P0044-B09
Pro $n=4$ prodit $W=1+r+r^4+r^9=2+2r$: hinc $W=2+2i$, quoties $k$ est formae $4\mu+1$, atque $W=2-2i$, quoties $k$ est formae $4\mu+3$.

% A544-P0044-B10
Pro $n=8$ habemus $W=1+r+r^4+r^9+r^{16}+r^{25}+r^{36}+r^{49}=2+4r+2r^4$ $=4r$.\quad Hinc erit

% A544-P0044-B11
\[
\begin{aligned}
W&=(1+i)\surd8,\quad\text{quoties }k\text{ est formae }8\mu+1\\
W&=(-1+i)\surd8,\quad\text{quoties }k\text{ est formae }8\mu+3\\
W&=(-1-i)\surd8,\quad\text{quoties }k\text{ est formae }8\mu+5\\
W&=(1-i)\surd8,\quad\text{quoties }k\text{ est formae }8\mu+7
\end{aligned}
\]

% A544-P0044-B12
Si $n$ est altior potestas binarii, statuamus $n=2^{2\varkappa}q$, ita ut $q$ sit vel $=1$ vel $=2$, atque $\varkappa$ maior quam $1$.\quad Hic ante omnia observari debet, si $\lambda$ sit integer quicunque per $2^{\varkappa-1}$ non divisibilis, fieri

% A544-P0044-B13
\[
\begin{aligned}
&r^{\lambda\lambda}+r^{(\lambda+2^\varkappa q)^2}+r^{(\lambda+2\pdot2^\varkappa q)^2}+r^{(\lambda+3\pdot2^\varkappa q)^2}\\&\qquad+\text{etc.}+r^{(\lambda+n-2^\varkappa q)^2}\\
&\quad=r^{\lambda\lambda}\{1+r^{2^{\varkappa+1}\lambda q}+r^{2\pdot2^{\varkappa+1}\lambda q}\\&\qquad+r^{3\pdot2^{\varkappa+1}\lambda q}+\text{etc.}+r^{(2n-2^{\varkappa+1}q)\lambda}\}\\&\quad=\frac{r^{\lambda\lambda}(1-r^{2\lambda n})}{1-r^{2^{\varkappa+1}\lambda q}}=0
\end{aligned}
\]

% A544-P0044-B14
Hinc facile perspicietur, fieri

% A544-P0044-B15
\[
\begin{aligned}
W=1+r^{2^{2\varkappa-2}}+r^{4\pdot2^{2\varkappa-2}}+r^{9\pdot2^{2\varkappa-2}}+\text{etc.}+r^{(n-2^{\varkappa-1})^2}
\end{aligned}
\]

\SourceLeaf{45}%
% A544-P0045-B02
Statuamus $r^{2^{2\varkappa-2}}=\rho$, eritque $\rho$ radix aequationis $x^{4q}-1=0$, et quidem $\rho=\cos\frac{k}{4q}360^0+i\sin\frac{k}{4q}360^0$; dein fiet

% A544-P0045-B03
\[
\begin{aligned}
W&=1+\rho+\rho^4+\rho^9+\text{etc.}+\rho^{(2^{\varkappa+1}q-1)^2}\\
&=2^{\varkappa-1}(1+\rho+\rho^4+\rho^9+\text{etc.}+\rho^{(4q-1)^2})
\end{aligned}
\]

% A544-P0045-B04
Sed summa seriei $1+\rho+\rho^4+\rho^9+\text{etc.}+\rho^{(4q-1)^2}$ per ea, quae de casibus $n=4$, $n=8$ explicavimus, determinatur, unde colligimus

% A544-P0045-B05
in casu eo, ubi $q=1$, sive ubi $n$ est potestas numeri $4$, fieri

% A544-P0045-B06
\[
\begin{aligned}
W&=(1+i)2^\varkappa=(1+i)\surd n,\quad\text{si fuerit }k\text{ formae }4\mu+1\\
W&=(1-i)2^\varkappa=(1-i)\surd n,\quad\text{si fuerit }k\text{ formae }4\mu+3
\end{aligned}
\]

% A544-P0045-B07
quae sunt ipsissimae formulae pro $n=4$ traditae;

% A544-P0045-B08
in casu eo autem, ubi $q=2$, sive ubi $n$ est potestas binarii cum exponente impari maiori quam $3$, fieri

% A544-P0045-B09
\[
\begin{aligned}
W&=(1+i)2^\varkappa\surd2=(1+i)\surd n,\quad\text{si fuerit }k\text{ formae }8\mu+1\\
W&=(-1+i)2^\varkappa\surd2=(-1+i)\surd n,\quad\text{si fuerit }k\text{ formae }8\mu+3\\
W&=(-1-i)2^\varkappa\surd2=(-1-i)\surd n,\quad\text{si fuerit }k\text{ formae }8\mu+5\\
W&=(1-i)2^\varkappa\surd2=(1-i)\surd n,\quad\text{si fuerit }k\text{ formae }8\mu+7
\end{aligned}
\]

% A544-P0045-B10
quae quoque prorsus conveniunt cum iis, quae pro $n=8$ tradidimus.

% A544-P0045-B11
\Article{24}

% A544-P0045-B12
Etiam hic operae pretium erit, rationem summae progressionis

% A544-P0045-B13
\[
\begin{aligned}
W'=1+r^h+r^{4h}+r^{9h}+\text{etc.}+r^{h(n-1)^2}
\end{aligned}
\]

% A544-P0045-B14
ad $W$ determinare, ubi $h$ integrum quemcunque imparem denotat.\quad Quum $W'$ oriatur ex $W$, mutando $k$ in $kh$, valor ipsius $W'$ perinde a forma numeri $kh$ pendebit, ut $W$ a forma ipsius $k$.\quad Statuamus $\frac{W'}{W}=l$, patetque

% A544-P0045-B15
I.\quad in casu eo, ubi $n=4$, vel altior potestas binarii cum exponente pari, fieri

% A544-P0045-B16
\[
\begin{aligned}
l&=1,\quad\text{si fuerit }h\text{ formae }4\mu+1\\
l&=-i,\quad\text{si fuerit }h\text{ formae }4\mu+3,\;\text{atque }k\text{ formae }4\mu+1\\
l&=+i,\quad\text{si fuerit }h\text{ formae }4\mu+3,\;\text{atque }k\text{ eiusdem formae.}
\end{aligned}
\]

\SourceLeaf{46}%
% A544-P0046-B02
II.\quad in casu eo, ubi $n=8$, vel altior potestas binarii cum exponente impari, fieri

% A544-P0046-B03
\[
\begin{aligned}
l&=1,\ \text{si fuerit}\  h\text{ formae }8\mu+1,\\
l&=-1,\ \text{si fuerit}\  h\text{ formae }8\mu+5,\\
l&=+i,\ \text{si fuerit vel }h\text{ formae }8\mu+3,\;\text{atque }k\text{ formae }4\mu+1,\\
&\phantom{=+i,}\text{vel }h\text{ formae }8\mu+7,\;\text{atque }k\text{ formae }4\mu+3,\\
l&=-i,\ \text{si fuerit vel }h\text{ formae }8\mu+3,\;\text{atque }k\text{ formae }4\mu+3,\\
&\phantom{=+i,}\text{vel }h\text{ formae }8\mu+7,\;\text{atque }k\text{ formae }4\mu+1.
\end{aligned}
\]

% A544-P0046-B04
Per praecc. determinatio summae $W$ pro iis casibus, ubi $n$ est numerus primus vel numeri primi potestas, complete perfecta est: superest itaque, ut eos quoque casus absolvamus, ubi $n$ e pluribus numeris primis compositus est, huc viam nobis sternet theorema sequens:

% A544-P0046-B05
\Article{25}

% A544-P0046-B06
\textsc{Theorema.}\quad \textit{Sit $n$ productum e duobus integris positivis inter se primis $a$, $b$,} \textit{statuaturque}

% A544-P0046-B07
\[
\begin{aligned}
P&=1+ r^{aa}+r^{4aa}+r^{9aa}+\text{etc.}+r^{(b-1)^2aa}\\
Q&=1+r^{bb}+r^{4bb}+r^{9bb}+\text{etc.}+r^{(a-1)^2bb}
\end{aligned}
\]

% A544-P0046-B08
\textit{Tum dico fore $W=PQ$.}

% A544-P0046-B09
\textit{Demonstr.}\quad Designet $\alpha$ indefinite numeros $0,1,2,3\dotsfour a-1$, $\beta$ indefinite numeros $0,1,2,3\dotsfour b-1$, $\nu$ indefinite numeros $0,1,2,3\dotsfour n-1$.\quad Tunc patet esse

% A544-P0046-B10
\[
\begin{aligned}
P=\Sigma r^{aa\beta\beta},\quad Q=\Sigma r^{bb\alpha\alpha},\quad W=\Sigma r^{\nu\nu}
\end{aligned}
\]

% A544-P0046-B11
Hinc erit $PQ=\Sigma r^{aa\beta\beta+bb\alpha\alpha}$, substituendo pro $\alpha$ et $\beta$ omnes valores, omnibus modis inter se combinatos; hinc porro propter $2ab\alpha\beta=2\alpha\beta n$, erit $PQ=\Sigma r^{(a\beta+b\alpha)^2}$. Sed nullo negotio perspicitur, singulos valores ipsius $a\beta+b\alpha$ inter se diversos esse, atque alicui valori ipsius $\nu$ aequales.\quad Hinc erit $PQ=\Sigma r^{\nu\nu}=W$.

% A544-P0046-B12
Ceterum notandum est, $r^{aa}$ esse radicem propriam aequationis $x^b-1=0$, atque $r^{bb}$ radicem propriam aequationis $x^a-1=0$.

\SourceLeaf{47}%
% A544-P0047-B02
\Article{26}

% A544-P0047-B03
Sit porro $n$ productum e tribus numeris inter se primis $a,b,c$, patetque, si statuatur $bc=b'$, etiam $a$ et $b'$ inter se primos fore; adeoque $W$ productum e duobus factoribus

% A544-P0047-B04
\[
\begin{aligned}
&1+r^{aa}+r^{4aa}+r^{9aa}+\text{etc.}+r^{(b'-1)^2aa}\\
&1+r^{b'b'}+r^{4b'b'}+r^{9b'b'}+\text{etc.}+r^{(a-1)^2b'b'}
\end{aligned}
\]

% A544-P0047-B05
Sed quum $r^{aa}$ sit radix propria aequationis $x^{bc}-1=0$, erit ipse factor prior productum ex

% A544-P0047-B06
\[
\begin{aligned}
&1+\rho^{bb}+\rho^{4bb}+\rho^{9bb}+\text{etc.}+\rho^{(c-1)^2bb}\\
&1+\rho^{cc}+\rho^{4cc}+\rho^{9cc}+\text{etc.}+\rho^{(b-1)^2cc}
\end{aligned}
\]

% A544-P0047-B07
si statuitur $r^{aa}=\rho$.\quad Hinc patet, $W$ esse productum e factoribus tribus

% A544-P0047-B08
\[
\begin{aligned}
&1+r^{bbcc}+r^{4bbcc}+r^{9bbcc}+\text{etc.}+r^{(a-1)^2bbcc}\\
&1+r^{aacc}+r^{4aacc}+r^{9aacc}+\text{etc.}+r^{(b-1)^2aacc}\\
&1+r^{aabb}+r^{4aabb}+r^{9aabb}+\text{etc.}+r^{(c-1)^2aabb}
\end{aligned}
\]

% A544-P0047-B09
ubi $r^{bbcc}$, $r^{aacc}$, $r^{aabb}$ erunt resp. radices propriae aequationum $x^a-1=0$, $x^b-1=0$, $x^c-1=0$.

% A544-P0047-B10
\Article{27}

% A544-P0047-B11
Hinc facile concluditur generaliter, si $n$ sit productum e factoribus quotcunque inter se primis $a,b,c$ etc., $W$ fieri productum e totidem factoribus, qui sint

% A544-P0047-B12
\[
\begin{aligned}
&1+r^{\frac{nn}{aa}}+r^{\frac{4nn}{aa}}+r^{\frac{9nn}{aa}}+\text{etc.}+r^{\frac{(a-1)^2nn}{aa}}\\
&1+r^{\frac{nn}{bb}}+r^{\frac{4nn}{bb}}+r^{\frac{9nn}{bb}}+\text{etc.}+r^{\frac{(b-1)^2nn}{bb}}\\
&1+r^{\frac{nn}{cc}}+r^{\frac{4nn}{cc}}+r^{\frac{9nn}{cc}}+\text{etc.}+r^{\frac{(c-1)^2nn}{cc}},\;\text{etc.}
\end{aligned}
\]

% A544-P0047-B13
ubi $r^{\frac{nn}{aa}}$, $r^{\frac{nn}{bb}}$, $r^{\frac{nn}{cc}}$ etc. erunt radices propriae aequationum $x^a-1=0$, $x^b-1=0$. $x^c-1=0$ etc.

% A544-P0047-B14
\Article{28}

% A544-P0047-B15
Ex his principiis transitus ad determinationem completam ipsius $W$ pro valore quocunque ipsius $n$ sponte iam obvius est.\quad Decomponatur scilicet $n$ in facto\SourceLeaf{48}%
% A544-P0048-B02
res $a,b,c$ etc. tales, qui sint vel numeri primi inaequales, vel potestates numerorum primorum inaequalium, statuatur $r^{\frac{nn}{aa}}=A$, $r^{\frac{nn}{bb}}=B$, $r^{\frac{nn}{cc}}=C$ etc. eruntque $A,B,C$ etc. radices propriae aequationum $x^a-1=0$, $x^b-1=0$, $x^c-1=0$ etc., atque $W$ productum e factoribus

% A544-P0048-B03
\[
\begin{aligned}
&1+A+A^4+A^9+\text{etc.}+A^{(a-1)^2}\\
&1+B+B^4+B^9+\text{etc.}+B^{(b-1)^2}\\
&1+C+C^4+C^9+\text{etc.}+C^{(c-1)^2}
\end{aligned}
\]

% A544-P0048-B04
Sed hi singuli factores per ea, quae in artt. $20,21,23$ docuimus, determinari poterunt, unde etiam valor producti innotescet.\quad Regulas pro determinandis illis factoribus hic in unum obtutum collegisse haud inutile erit.\quad Quum radix $A$ fiat $=\frac{kn}{a}\pdot\frac{360^0}{a}$, aggregatum $1+A+A^4+A^9+\text{etc.}+A^{(a-1)^2}$, quod per $L$ denotabimus, perinde per numerum $\frac{kn}{a}$ determinabitur, ut in disquisitione nostra generali $W$ per $k$.\quad Duodecim iam casus sunt distinguendi.

% A544-P0048-B05
I.\quad Si $a$ est numerus primus formae $4\mu+1$, puta $=p$, vel potestas talis numeri primi cum exponente impari, simulque $\frac{kn}{a}$ residuum quadraticum ipsius $p$, erit $L=+\surd a$.

% A544-P0048-B06
II.\quad Si manentibus reliquis $\frac{kn}{a}$ est non-residuum quadraticum ipsius $p$, erit $L=-\surd a$.

% A544-P0048-B07
III.\quad Si $a$ est numerus primus formae $4\mu+3$, puta $=p$, vel potestas talis numeri primi cum exponente impari, simulque $\frac{kn}{a}$ residuum quadraticum ipsius $p$, erit $L=+i\surd a$.

% A544-P0048-B08
IV.\quad Si, manentibus reliquis ut in III, $\frac{kn}{a}$ est non-residuum quadraticum ipsius $p$, erit $L=-i\surd a$.

% A544-P0048-B09
V.\quad Si $a$ est quadratum, altiorve potestas numeri primi (imparis) cum exponente pari, erit $L=+\surd a$.

% A544-P0048-B10
VI.\quad Si $a=2$, erit $L=0$.

% A544-P0048-B11
VII.\quad Si $a=4$, altiorve potestas binarii cum exponente pari, simulque $\frac{kn}{a}$ formae $4\mu+1$, erit $L=(1+i)\surd a$.

% A544-P0048-B12
VIII.\quad Si, manentibus reliquis ut in VII, $\frac{kn}{a}$ est formae $4\mu+3$, erit $L=(1-i)\surd a$.

% A544-P0048-B13
IX.\quad Si $a=8$, altiorve potestas binarii cum exponente impari, simulque $\frac{kn}{a}$ formae $8\mu+1$, erit $L=(1+i)\surd a$.

\SourceLeaf{49}%
% A544-P0049-B02
X.\quad Si, manentibus reliquis ut in IX, $\frac{kn}{a}$ est formae $8\mu+3$, erit $L=(-1+i)\surd a$.

% A544-P0049-B03
XI.\quad Si manentibus reliquis $\frac{kn}{a}$ est formae $8\mu+5$, erit $L=(-1-i)\surd a$.

% A544-P0049-B04
XII.\quad Si manentibus reliquis $\frac{kn}{a}$ est formae $8\mu+7$, erit $L=(1-i)\surd a$.

% A544-P0049-B05
\Article{29}

% A544-P0049-B06
Sit exempli caussa $n=2520=8\pdot9\pdot5\pdot7$, atque $k=13$.\quad Hic erit

% A544-P0049-B07
{\small
\[
\begin{aligned}
&\text{pro }a=8,\;\text{per casum XII,}\quad L=(1-i)\surd8\\
&\text{pro factore }9,\;\text{per casum V, summa respondens erit }=\surd9\\
&\text{pro factore }5,\;\text{per casum II, summa respondens erit }=-\surd5\\
&\text{pro factore }7,\;\text{per casum III, summa respondens erit }=+i\surd7\mathord{.}
\end{aligned}
\]
}

% A544-P0049-B08
Hinc fit $W=(1-i)\pdot(-i)\pdot\surd2520=(-1-i)\surd2520$.

% A544-P0049-B09
Sit pro eodem valore ipsius $n$, $k=1$: tunc respondebit

% A544-P0049-B10
\[
\begin{aligned}
&\text{factori }8\;\text{summa }(-1+i)\surd8\\
&\text{factori }9\;\text{summa }\surd9\\
&\text{factori }5\;\text{summa }\surd5\\
&\text{factori }7\;\text{summa }-i\surd7
\end{aligned}
\]

% A544-P0049-B11
Hinc conflatur productum $W=(1+i)\surd2520$.

% A544-P0049-B12
\Article{30}

% A544-P0049-B13
Methodus alia, summam $W$ generaliter determinandi, petitur ex iis, quae in artt. $22.\ 24$ exposita sunt.\quad Statuamus $\cos\omega+i\sin\omega=\rho$, atque

% A544-P0049-B14
\[
\begin{aligned}
\rho^{\frac{nn}{aa}}=\alpha,\quad\rho^{\frac{nn}{bb}}=\beta,\quad\rho^{\frac{nn}{cc}}=\gamma\;\text{etc.}
\end{aligned}
\]

% A544-P0049-B15
ita ut habeatur $r=\rho^k$, $A=\alpha^k$, $B=\beta^k$, $C=\gamma^k$ etc.\quad Tunc erit

% A544-P0049-B16
\[
\begin{aligned}
1+\rho+\rho^4+\rho^9+\text{etc.}+\rho^{(n-1)^2}
\end{aligned}
\]

% A544-P0049-B17
productum e factoribus

% A544-P0049-B18
\[
\begin{aligned}
&1+\alpha+\alpha^4+\alpha^9+\text{etc.}+\alpha^{(a-1)^2}\\
&1+\beta+\beta^4+\beta^9+\text{etc.}+\beta^{(b-1)^2}\\
&1+\gamma+\gamma^4+\gamma^9+\text{etc.}+\gamma^{(c-1)^2},\;\text{etc.}
\end{aligned}
\]

\SourceLeaf{50}%
% A544-P0050-B02
adeoque $W$ productum e factoribus

% A544-P0050-B03
\[
\begin{aligned}
w&=1+\rho+\rho^4+\rho^9+\text{etc.}+\rho^{(n-1)^2}\\[2pt]
\mathfrak A&=\frac{1+A+A^4+A^9+\text{etc.}+A^{(a-1)^2}}{1+\alpha+\alpha^4+\alpha^9+\text{etc.}+\alpha^{(a-1)^2}}\\[2pt]
\mathfrak B&=\frac{1+B+B^4+B^9+\text{etc.}+B^{(b-1)^2}}{1+\beta+\beta^4+\beta^9+\text{etc.}+\beta^{(b-1)^2}}\\[2pt]
\mathfrak C&=\frac{1+C+C^4+C^9+\text{etc.}+C^{(c-1)^2}}{1+\gamma+\gamma^4+\gamma^9+\text{etc.}+\gamma^{(c-1)^2}}
\end{aligned}
\]

% A544-P0050-B04
Iam factor primus $w$ determinatus est per disquisitiones supra traditas (art. $19$); factores reliqui vero $\mathfrak A,\mathfrak B,\mathfrak C$ etc. prodeunt per formulas artt. $22.\ 24$, quas ut omnia iuncta habeantur, hic denuo colligimus\begingroup\renewcommand{\thefootnote}{*)}\footnote{Manifesto, quae illic erant $k$ et $h$, hic erunt $\frac{n}{a}$ et $k$ respectu factoris secundi, $\frac{n}{b}$ et $k$ respectu factoris tertii etc.}\endgroup.\quad Duodecim casus hic sunt distinguendi, scilicet

% A544-P0050-B05
I.\quad Si $a$ est numerus primus (impar) $=p$, vel talis numeri potestas cum exponente impari, atque $k$ residuum quadraticum ipsius $p$, erit factor respondens $\mathfrak A=+1$.

% A544-P0050-B06
II.\quad Si manentibus reliquis $k$ est non-residuum quadraticum ipsius $p$, erit $\mathfrak A=-1$.

% A544-P0050-B07
III.\quad Si $a$ est quadratum numeri primi imparis, altiorve eius potestas cum exponente pari, erit $\mathfrak A=+1$.

% A544-P0050-B08
IV.\quad Si $a$ est $=4$, aut altior binarii potestas cum exponente pari, simulque $k$ formae $4\mu+1$, erit $\mathfrak A=+1$.

% A544-P0050-B09
V.\quad Si, manentibus reliquis ut in IV, $k$ est formae $4\mu+3$, atque $\frac{n}{a}$ formae $4\mu+1$, erit $\mathfrak A=-i$.

% A544-P0050-B10
VI.\quad Si, manentibus reliquis ut in IV, $k$ est formae $4\mu+3$, atque $\frac{n}{a}$ formae $4\mu+3$, erit $\mathfrak A=+i$.

% A544-P0050-B11
VII.\quad Si $a$ est $=8$, aut altior binarii potestas cum exponente impari, atque $k$ formae $8\mu+1$, erit $\mathfrak A=+1$.

% A544-P0050-B12
VIII.\quad Si, manentibus reliquis ut in VII, $k$ est formae $8\mu+5$, erit $\mathfrak A=-1$.

% A544-P0050-B13
IX.\quad Si, manentibus reliquis ut in VII, $k$ est formae $8\mu+3$, atque $\frac{n}{a}$ formae $4\mu+1$, erit $\mathfrak A=+i$.

